\documentclass[11pt,a4paper]{article}

\usepackage[T1]{fontenc}
\usepackage[utf8]{inputenc}
\usepackage{lmodern}
\usepackage{amsmath,amssymb}
\usepackage{geometry}
\usepackage{hyperref}

\geometry{margin=1in}

\title{Theta, Vector, and Unified-Action Constraints in TRM/TQM V3.0}
\author{
Fabrice Wieser\\
Independent Researcher\\
\url{https://github.com/MagusDraconis/TRM_Cosmology}
}
\date{June 2026\\
V3.0 Review Baseline\\
Release tag: \texttt{trm-v3.0-review-release}}

\begin{document}
\maketitle

\begin{abstract}
This paper isolates the TRM/TQM V3.0 theta-observable path, vector frame-dragging candidate sector, and unified effective-action roadmap \cite{theta_theorem_path,vector_extension,unified_action_roadmap,trm_status_v3}. The objective is to state what is currently tested-effective, what is structurally constrained, and what remains open at theorem level. The chain \(\Theta \rightarrow O_5 \rightarrow \lambda_\Theta \rightarrow g_{\mathrm{obs}}\) is strongly supported by guard tests; the vector sector is strongly hardened in weak-field candidate form through FD01-FD20; and unified-action guards UF01-UF09 preserve sector limits and bounded cross-couplings. These results remain below theorem-level first-principles closure and do not claim GR replacement.
\end{abstract}

\section{Introduction}

TRM/TQM V3.0 uses sector-wise hardening and cross-sector consistency guards to constrain model freedom while keeping claim boundaries explicit \cite{trm_status_v3}. This paper focuses on the three tracks that define current cross-sector structure:
\begin{enumerate}
  \item the nonlocal theta-observable chain,
  \item the vector rotational/frame-dragging candidate sector,
  \item the unified effective-action roadmap that links scalar, vector, and theta sectors.
\end{enumerate}

\section{Theta-Observable Path}

\subsection{Core chain}

\[
\Theta \rightarrow O_5 \rightarrow \lambda_\Theta \rightarrow g_{\mathrm{obs}}.
\]

\subsection{Current test support}

The theta path is currently hardened by TO01-TO28, TQK01-TQK04, LC01-LC08, and TOL01-TOL04 \cite{theta_theorem_path,trm_status_v3}. At this stage, the chain is treated as tested-effective with strong guard-level support for operator behavior, holdout discipline, anti-proxy safeguards, and lattice-energy consistency.

\subsection{Claim boundary}

\begin{itemize}
  \item Supported: strongly constrained theorem path at tested-effective level.
  \item Not claimed: theorem-level microscopic closure of the full \(\Theta \rightarrow g_{\mathrm{obs}}\) map.
\end{itemize}

\section{Vector Frame-Dragging Candidate Sector}

\subsection{Core structure}

The vector sector is based on:
\[
\vec A_T,\qquad \vec B_T=\nabla\times\vec A_T,
\]
with weak-field frame-dragging candidate interpretation \cite{vector_extension,lense_thirring}.

\subsection{Current hardening status}

FD01-FD20 currently support \cite{vector_extension,trm_status_v3}:
\begin{enumerate}
  \item rotational source generation and spin-sign reversal behavior,
  \item weak-field scaling shape and null/spin-zero fallback consistency,
  \item derived non-fitted effective \(k_T\) normalization discipline,
  \item ablation/holdout/unit-scaling robustness,
  \item controlled systematic-bias auditing under no-refit constraints.
\end{enumerate}

\subsection{Claim boundary}

\begin{itemize}
  \item Supported: tested-effective weak-field frame-dragging candidate sector.
  \item Not claimed: theorem-level first-principles closure or full quantitative GR-equivalence.
\end{itemize}

\section{Unified Effective-Action Roadmap}

\subsection{Canonical roadmap form}

\[
S_{\mathrm{eff}}[T,\vec A_T,\Theta]
=
S_T[T]+S_A[\vec A_T]+S_\Theta[\Theta]+S_{\mathrm{int}}[T,\vec A_T,\Theta].
\]

\subsection{Current guard status}

UF01-UF09 currently support \cite{unified_action_roadmap,trm_status_v3}:
\begin{enumerate}
  \item scalar/vector/theta limit recovery,
  \item zero-coupling additive decomposition,
  \item bounded small cross-couplings,
  \item global cross-coupling identifiability without per-group refit dependence,
  \item preservation of MC/FD/TO guard behavior under small couplings.
\end{enumerate}

\subsection{Claim boundary}

\begin{itemize}
  \item Supported: candidate-level unified effective-action roadmap with guard preservation.
  \item Not claimed: theorem-level unification.
\end{itemize}

\section{Cross-Sector Falsification Criteria}

The current theta/vector/unified interpretation should be narrowed or rejected if:
\begin{enumerate}
  \item theta-chain anti-proxy or holdout guarantees collapse under independent tests;
  \item vector weak-field guard behavior fails under robust ablation/benchmark stress;
  \item unified-action interaction terms break sector-limit recovery or MC/FD/TO guard preservation;
  \item claimed cross-coupling structure requires hidden per-dataset re-normalization.
\end{enumerate}

\section{Remaining First-Principles Gaps}

Open items remain explicit \cite{theta_theorem_path,unified_action_roadmap,vector_extension}:
\begin{itemize}
  \item theorem-level microscopic closure for the theta chain,
  \item theorem-level first-principles closure for vector-sector normalization and benchmark mapping,
  \item theorem-level closure of unified cross-sector interaction structure.
\end{itemize}

\section{Conclusion}

In V3.0, theta, vector, and unified-action blocks are strongly constrained at tested-effective level and form a coherent candidate multi-sector framework under explicit guard discipline. The status is review-ready and falsifiable, but it is not theorem-level complete and does not claim replacement of General Relativity.

\bibliographystyle{plain}
\bibliography{references}

\end{document}
